% 76-08(76쪽) — 두 함수 $y=\log_3 x$ · $y=x$ 의 그래프와 $x_1$·$x_2$·$x_3$ 을 잇는 계단 모양 안내 점선.
% 스캔 화질이 낮아 TikZ 로 다시 그림. 원문 그림은 눈금도 비례도 없는 개형이라(오른쪽으로 갈수록 가로축이 줄어든 왜곡이라
% 그려진 곡선이 실제 $\log_3 x$ 가 아니다) 구조($x_1$ 은 곡선의 $x$절편 · 세로 점선으로 $y=x$ 까지 올라간 높이에서
% 가로로 그으면 곡선과 만나는 곳이 다음 $x$)와 배치 비율만 원문을 따랐다.
% 눈금·좌표값은 적지 않는다 — $x_1+x_2+x_3$ 의 값(답)이 그림에서 읽히지 않게 한다.
\begin{tikzpicture}[line width=0.6pt, x=0.3cm, y=0.3cm]
  \draw[->, >=stealth, line width=0.7pt] (-1.2,0) -- (12.6,0) node[below=1pt, font=\small] {$x$};
  \draw[->, >=stealth, line width=0.7pt] (0,-1.7) -- (0,8.9) node[left=1pt, font=\small] {$y$};
  \draw[line width=0.5pt] (-1.0,-1.0) -- (7.1,7.1);
  \draw[line width=0.6pt, smooth] plot coordinates
    {(0.60,-1.05) (0.70,-0.65) (0.80,-0.38) (0.90,-0.15) (1.00,0) (1.10,0.08)
     (1.30,0.18) (1.60,0.31) (2.00,0.49) (2.60,0.76) (3.20,1.00) (4.20,1.39)
     (5.50,1.90) (7.00,2.49) (8.80,3.20) (10.20,3.75) (11.40,4.22)};
  \draw[densely dotted, line width=0.4pt]
    (1,0) -- (1,1) (0,1) -- (3.2,1) (3.2,0) -- (3.2,3.2) (0,3.2) -- (8.8,3.2) (8.8,0) -- (8.8,3.2);
  \node[below=1pt, xshift=3.5pt, font=\small] at (1,0) {$x_1$};
  \node[below=1pt, font=\small] at (3.2,0) {$x_2$};
  \node[below=1pt, font=\small] at (8.8,0) {$x_3$};
  \node[above left=-2.5pt and -1pt, font=\small] at (0,0) {$\mathrm{O}$};
  \node[anchor=west, font=\small] at (7.0,8.0) {$y=x$};
  \node[anchor=west, font=\small] at (6.4,4.5) {$y=\log_3 x$};
\end{tikzpicture}
